%% kapa: exam=1 round=37 year=2026 period=2 form=A subject=appraisal_law subject_ko=감정평가관계법규 seq=19 number=19
%% source: Q-Net 감정평가사 기출문제 (cst00309) artlSeq=5258924 file=raw/1008/5258924/2026년 제37회 감정평가사 1차 2교시.pdf
%% conversion: kordoc 4.24.51 -> tools/gen1.py (auto); review=proofread-sonnet; flags=none
\begin{question}{19}
감정평가 및 감정평가사에 관한 법령상 감정평가법인등에 대한 과징금 부과에 관한 설명으로 옳은 것은?
\begin{choices}
  \item 과징금 부과에 대하여 이의신청이 있는 경우, 이의신청에 대한 결정에 이의가 있는 자는 ｢행정심판법｣에 따른 행정심판을 청구할 수 없다.
  \item 과징금납부의무자가 과징금 납부기한을 연장받거나 분할납부를 하려면 납부기한 7일 전까지 국토교통부장관에게 신청하여야 한다.
  \item 과징금 납부기한의 연장은 납부기한의 다음 날부터 1년을 초과할 수 없다.
  \item 과징금을 분할납부하게 하는 경우 각 분할된 납부기한 간의 간격은 3개월 이내로 하며, 분할 횟수는 3회 이내로 한다.
  \item 과징금납부의무자가 납부기한까지 과징금을 납부하지 아니하여 독촉장을 발부받은 경우 체납된 과징금의 납부기한은 독촉장 발부일부터 15일 이내로 한다.
\end{choices}
\end{question}
\answer{3}
